\section{Resolución de sistemas lineales}
\subsection{Introducción a los sistemas lineales}

Sea $A \in \mathbb{K}^{n \times n}$ y $b \in \mathbb{K}^n$. Un sistema lineal es la búsqueda de $x \in \mathbb{K}^n$ tal que $S \equiv Ax=b$. En este apartado se estudiarán diferentes métodos para resolver sistemas lineales, así como sus propiedades y análisis de convergencia.\\

Es fácil ver que si $A$ es invertible, entonces el sistema tiene una única solución dada por $x=A^{-1}b$. Sin embargo, calcular la inversa de una matriz puede ser computacionalmente costoso y numéricamente inestable. Por ello, se han desarrollado diversos métodos para resolver sistemas lineales sin necesidad de calcular la inversa de la matriz.\\

\begin{teorema}[Metodo de Cramer]
    Sea el sistema $S \equiv Ax=b$ con $A \in \mathbb{K}^{n \times n}$ invertible. Entonces, si denotamos:
    \[
        A_i = \begin{bmatrix}
            a_{11} & \cdots & a_{1,i-1} & b_1 & a_{1,i+1} & \cdots & a_{1n} \\
            a_{21} & \cdots & a_{2,i-1} & b_2 & a_{2,i+1} & \cdots & a_{2n} \\
            \vdots  &  & \vdots  & \vdots  & \vdots  &  & \vdots  \\
            a_{n1} & \cdots & a_{n,i-1} & b_n & a_{n,i+1} & \cdots & a_{nn}
        \end{bmatrix} =
        \left( a_1, \ldots, a_{i-1}, b, a_{i+1}, \ldots, a_n \right)
    \]
    Dónde $a_j$ es la $j$-ésima columna de $A$, entonces la solución del sistema es:
    \[
        x_i = \frac{\det(A_i)}{\det(A)} \quad \forall i=1, \ldots, n
    \]
\end{teorema}

Sin embargo, el método de Cramer es computacionalmente ineficiente para matrices grandes, ya que el cálculo de determinantes tiene una complejidad factorial. Por ello, se utilizan otros métodos como el método de Gauss, que es mucho más eficiente y estable para resolver sistemas lineales, se pueden apreciar los tiempos necesarios en la siguiente tabla:

\[
\begin{array}{|c|c|c|c|}
    \hline
    n & 10 & 100 & 10^3 \\
    \hline
    \text{Cramer} & 10^8 & 10^{162} & 10^{2573} \\
    \hline
    \text{Gauss} & 800 & 10^6 & 10^9 \\
    \hline
\end{array}
\]

Siendo $n$ el tamaño de la matriz $A$, el metodo de Cramer tiene una complejidad de $O(n^2 \cdot n!)$ mientras que el método de Gauss tiene una complejidad de $O(n^3)$.\\

En los siguientes apartados se estudiarán diferentes métodos para resolver sistemas lineales, así como sus propiedades y análisis de convergencia.\\

Un caso particular de sistemas lineales es el caso de sistemas diagonales, que tienen esta forma:

\[
    \begin{pmatrix}
        a_{11} & 0 & \cdots & 0 \\ 0 & a_{22} & \cdots & 0 \\ \vdots & \vdots & \ddots & \vdots \\ 0 & 0 & \cdots & a_{nn} 
    \end{pmatrix}
    \cdot x = b
\]

De lo que se deduce trivialmente que la solución es $x_i = \frac{b_i}{a_{ii}}$ para $i=1, \ldots, n$.\\

Estos sistemas diagonales tiene  una complejidad de $O(n)$, ya que se necesitan $n$ operaciones para calcular la solución.\\

Otro sistema particular es el sistema triangular superior, que tiene esta forma:

\[ 
\begin{pmatrix} 
    a_{11} & a_{12} & \cdots & a_{1n} \\ 0 & a_{22} & \cdots & a_{2n} \\ \vdots & \vdots & \ddots & \vdots \\ 0 & 0 & \cdots & a_{nn} 
\end{pmatrix} 
\cdot x = b 
\] 

En este caso, se puede resolver el sistema mediante sustitución hacia atrás, comenzando por la última ecuación y avanzando hacia la primera, es decir:

\[
    x_i = 
    \frac{1}{a_{ii}} \left( b_i - \sum_{j=i+1}^n a_{ij} x_j \right) 
    \quad \forall i \in \{n, n-1, \ldots, 1\} \
\]

Analogamente, un sistema triangular inferior tiene esta forma: 
\[ 
    \begin{pmatrix} a_{11} & 0 & \cdots & 0 \\ a_{21} & a_{22} & \cdots & 0 \\ \vdots & \vdots & \ddots & \vdots \\ a_{n1} & a_{n2} & \cdots & a_{nn} 
    \end{pmatrix} 
    \cdot x = b 
\] 

Y se puede resolver mediante sustitución hacia adelante, comenzando por la primera ecuación y avanzando hacia la última, es decir: 
\[ 
x_i = \frac{1}{a_{ii}} \left( b_i - \sum_{j=1}^{i-1} a_{ij} x_j \right) 
\quad \forall i \in \{1, 2, \ldots, n\}
\]

Estos sistemas triangulares tienen una complejidad de $O(n^2)$, ya que se necesitan $n$ operaciones para calcular cada $x_i$, y hay $n$ variables.\\

Ahora veamos algunos metodos de resolución de sistemas lineales más generales.

\begin{proposición}[Método de eliminación Gaussiana]
    El método de eliminación Gaussiana es un algoritmo para resolver sistemas lineales que consiste en transformar la matriz $A$ en una matriz triangular superior mediante operaciones elementales de fila, es decir:

    \[
    Ax = b \longrightarrow
    Tx = b'
    \]

    Para pasar a una matriz triangular superior $T$, usaremos una matriz $E$ cuyo determinante no sea nula, de tal forma que tengamos el siguiente sistema equivalente:

    \[
    \underbrace{E \cdot A}_{T} x = \underbrace{E \cdot b}_{b'}
    \]
    
    Para ver quien es esta matriz $E$, se pueden usar matrices de permutación para intercambiar filas, matrices de escala para multiplicar una fila por un escalar distinto de cero, y matrices de eliminación para sumar un múltiplo de una fila a otra fila. Denotemos:

    \[
    E = E_k \cdot E_{k-1} \cdots E_1
    \]

    Siendo $E_i$ una matriz elemental que representa una de las operaciones mencionadas anteriormente, y siendo:

    \[
        E_i \cdot A_{i-1} = A_i \, \quad b_i = E_i \cdot b_{i-1} \longrightarrow
        A_i x = b_i
    \]

    Y por tanto el algoritmo se basara en llegar a un $k$ que cumpla:

    \[
        A_k = E_k \cdot E_{k-1} \cdots E_1 \cdot A = E \cdot A = T
    \]

    Ahora veamos como encontrar una matriz $E_i$ que nos permita hacer cada una de las operaciones elementales descritas anteriormente:
    \begin{itemize}
        \item \textbf{Multiplicar una fila por un escalar distinto de cero:} Para multiplicar la fila $i$ por un escalar $\lambda \neq 0$, se puede usar la matriz de escala $E_i$ que tiene la forma:
        \[
            E_i = \begin{pmatrix}
                1 & \cdots & 0 \\ 
                0 & \ddots & \cdots & 0 
                \\ \vdots & 0 & \lambda & 0 
                \\ 0 & 0 & \cdots & 1 
            \end{pmatrix}
        \]
        Donde el elemento en la posición $(i,i)$ es $\lambda$ y el resto de los elementos son $1$ en la diagonal y $0$ en el resto de la matriz, y transforma $A_{i-1}$ en $A_i$ multiplicando la fila $i$ por $\lambda$:
        \[
            E_i \cdot A_{i-1} =
            \begin{pmatrix}
                1 & 0 & \cdots & 0 \\ 
                0 & \ddots & \cdots & 0 
                \\ \vdots & 0 & \lambda & 0 
                \\ 0 & 0 & \cdots & 1
            \end{pmatrix}
            \cdot \begin{pmatrix}
                a_i^T \\ \vdots \\ a_i^T \\ \vdots \\ a_n^T 
            \end{pmatrix} =
            \begin{pmatrix}
                a_i^T \\ \vdots \\ \lambda a_i^T \\ \vdots \\ a_n^T
            \end{pmatrix}
        \]
        \item \textbf{Intercambiar dos filas:} Para intercambiar las filas $i$ y $j$, se puede usar la matriz de permutación $E_i$ que tiene la forma:
        \[
            E_i = \begin{pmatrix}
                1 & \cdots & 0 & \cdots & 0 \\ 
                \vdots & \ddots & \vdots & \ddots & \vdots 
                \\ 0 & \cdots & 0 & \cdots & 1 
                \\ \vdots & \ddots & \vdots & \ddots & \vdots 
                \\ 0 & \cdots & 1 & \cdots & 0
            \end{pmatrix}
        \]
        Donde el elemento en la posición $(i,j)$ es $1$, el elemento en la posición $(j,i)$ es $1$, y el resto de los elementos son $1$ en la diagonal y $0$ en el resto de la matriz, y transforma $A_{i-1}$ en $A_i$ intercambiando las filas $i$ y $j$:
        \[
            E_i \cdot A_{i-1} =
            \begin{pmatrix}
                1 & 0 & \cdots & 0 & \cdots & 0 \\ 
                0 & \ddots & \vdots & \ddots & \vdots & \vdots 
                \\ \vdots & 0 & 0 & \cdots & 1 & 0 
                \\ 0 & \cdots & 1 & \cdots & 0 & 0 
                \\ \vdots & \ddots & \vdots & \ddots & \vdots & \vdots 
                \\ 0 & \cdots & 0 & \cdots & 0 & 1
            \end{pmatrix}
            \cdot
            \begin{pmatrix}
                a_1^T \\ \vdots \\ a_i^T \\ \vdots \\ a_j^T \\ \vdots \\ a_n^T
            \end{pmatrix} =
            \begin{pmatrix}
                a_1^T \\ \vdots \\ a_j^T \\ \vdots \\ a_i^T \\ \vdots \\ a_n^T
            \end{pmatrix}
        \]
        \item \textbf{Suma de dos filas:} Para sumar una fila $i$ a otra fila $j$, se puede usar la matriz de eliminación $E_i$ que tiene la forma:
        \[
            E_i = \begin{pmatrix}
                1 & \cdots & 0 & \cdots & 0 \\ 
                \vdots & \ddots & \vdots & \ddots & \vdots 
                \\ 0 & \cdots & 1 & \cdots & 1
                \\ \vdots & \ddots & \vdots & \ddots & \vdots 
                \\ 0 & \cdots & 0 & \cdots & 1 
            \end{pmatrix}
        \]
        Donde el elemento en la posición $(j,i)$ es $1$, el elemento en la posición $(j,j)$ es $1$, y el resto de los elementos son $1$ en la diagonal y $0$ en el resto de la matriz, y transforma $A_{i-1}$ en $A_i$ sumando la fila $i$ a la fila $j$:
        \[
            E_i \cdot A_{i-1} =
            \begin{pmatrix}
                1 & 0 & \cdots & 0 & \cdots & 0 \\ 
                0 & \ddots & \vdots & \ddots & \vdots & \vdots 
                \\ 0 & \cdots & 1 & \cdots & 1 & 0 
                \\ 0 & \cdots & 0 & \cdots & 1 & 0 
                \\ \vdots & \ddots & \vdots & \ddots & \vdots & \vdots 
                \\ 0 & \cdots & 0 & \cdots & 0 & 1
            \end{pmatrix}
            \cdot
            \begin{pmatrix}
                a_1^T \\ \vdots \\ a_i^T \\ \vdots \\ a_j^T \\ \vdots \\ a_n^T
            \end{pmatrix} =
            \begin{pmatrix}
                a_1^T \\ \vdots \\ a_i^T + a_j^T \\ \vdots \\ a_j^T \\ \vdots \\ a_n^T
            \end{pmatrix}
        \]
    \end{itemize}
    
    Tras comprobar esto, veamos como aplicarlo a un sistema lineal, suponiendo que $\det(A) \neq 0$, entonces el sistema tiene una única solución, y podemos aplicar el método de eliminación Gaussiana para resolverlo.\\

    \begin{enumerate}
        \item Como $\det(A) \neq 0$, entonces $\exists i \in \{1, \ldots, n\}$ tal que $a_{i1} \neq 0$, entonces podemos usar la matriz de permutación $E_1$ para intercambiar la fila $i$ con la fila $1$, y obtener una nueva matriz $A_1$ con un elemento $a$ no nulo en la posición $(1,1)$:
        \[
            E_1 \cdot A = A_1 =
            \left(
            \begin{array}{c|c}
                a & \beta^T \\ \hline
                \alpha & A_{n-1}
            \end{array}
            \right)
        \]
        \item Despues, multiplicamos la fila $1$ por $\frac{1}{a}$ usando la matriz de escala $E_2$, y obtenemos una nueva matriz $A_2$ con un elemento $1$ en la posición $(1,1)$:
        \[
            E_2 \cdot A_1 = A_2 =
            \left(
            \begin{array}{c|c}
                1 & \frac{\beta^T}{a} \\ \hline
                \alpha & A_{n-1}
            \end{array}
            \right)
        \]
        \item Finalmente, sumamos un múltiplo de la fila $1$ a cada una de las filas $2, \ldots, n$ segun los elementos de $\alpha$ usando las matrices de eliminación $E_{l1}, \ldots, E_{ln}$, las cuales componen $E_3$, y obtenemos una nueva matriz $A_3$ con ceros en la primera columna excepto en la posición $(1,1)$:
        \[
            E_n \cdots E_3 \cdot A_2 = A_n =
            \left(
            \begin{array}{c|c}
                1 & \frac{\beta^T}{a} \\ \hline
                0 & A_{n-1} - \frac{\alpha \cdot \beta^T}{a}
            \end{array}
            \right)
        \]
        \item Si ahora denotamos $B_{n-1} = A_{n-1} - \frac{\alpha \cdot \beta^T}{a}$, entonces podemos repetir este proceso con la matriz $B_{n-1}$, y así sucesivamente hasta obtener una matriz triangular superior $T$:
        \[
            E_k \cdots E_3 \cdot A_2 = T
        \]
    \end{enumerate}
    Tras aplicar estos pasos, obtenemos una matriz triangular superior $T$ y un vector $b'$ tal que el sistema $Tx = b'$ es equivalente al sistema original $Ax = b$. Luego, podemos resolver el sistema triangular superior $Tx = b'$ mediante sustitución hacia atrás para obtener la solución del sistema original.\\
\end{proposición}

